Cette annexe décrit la structure organisationnelle de RUHENNA. Il s’agit d’une activité artisanale à taille humaine, centrée autour de sa fondatrice. L’objectif est de clarifier le contexte : la plateforme développée doit être administrable sans équipe technique dédiée, ce qui influence directement les choix d’architecture (simplicité, monorepo Next.js, back-office intégré).

\section{Organisation générale}
\begin{itemize}
    \item \textbf{Direction et production} : pilotées par Sumeyra (vision, fabrication, prestations).
    \item \textbf{Communication} : réseaux sociaux, contenu (photos, articles), interaction avec la clientèle.
    \item \textbf{Relation client} : demandes, devis, suivi des rendez-vous.
\end{itemize}

\section{Organigramme simplifié}
\begin{figure}[H]
    \centering
    \begin{adjustbox}{max width=\textwidth}
    \begin{tikzpicture}[
        node distance=8mm,
        box/.style={draw, rounded corners, align=center, inner sep=7pt, text width=0.55\textwidth},
        smallbox/.style={draw, rounded corners, align=center, inner sep=6pt, text width=0.26\textwidth},
        arrow/.style={-{Stealth[length=2.2mm]}, thick}
    ]
    \node[box] (sumeyra) {\textbf{Sumeyra Solak}\\\small Fondatrice RUHENNA\\Pilotage global (production, prestations, communication)};
    \node[smallbox, below left=of sumeyra] (prod) {Fabrication\\\small cônes, préparation};
    \node[smallbox, below=of sumeyra] (events) {Prestations\\\small événements, mariages};
    \node[smallbox, below right=of sumeyra] (com) {Communication\\\small galerie, blog, réseaux};
    \node[smallbox, below=of events] (client) {Relation client\\\small demandes, RDV, suivi};
    \node[smallbox, right=25mm of client] (stage) {Stagiaire\\\small conception + développement\\plateforme web};
    \draw[arrow] (sumeyra) -- (prod);
    \draw[arrow] (sumeyra) -- (events);
    \draw[arrow] (sumeyra) -- (com);
    \draw[arrow] (sumeyra) -- (client);
    \draw[arrow] (stage) -- (client);
    \draw[arrow] (stage) -- (com);
    \end{tikzpicture}
    \end{adjustbox}
    \caption{Organigramme simplifié de RUHENNA et positionnement du stage.}
    \label{fig:org_ruhenna}
\end{figure}

\section{Conséquences sur le produit}
Cette structure implique des exigences fortes :
\begin{itemize}
    \item \textbf{Autonomie} : le back-office doit être intuitif (ajout d’images, tri, publication).
    \item \textbf{Temps limité} : les workflows doivent éviter les tâches répétitives (centralisation messages/rdv).
    \item \textbf{Évolutivité} : possibilité d’ajouter un e-commerce complet plus tard sans réécrire le socle.
\end{itemize}
